\section{Classical Analysis Gaps: Full Scope}
\label{sec:classical-analysis-gaps}

\begin{layman}
The blueprint dependency graph for this project carries eight nodes
that show up with a red border: Radó's triangulation theorem, Tietze
normal form, the polygonal model, the de Rham theorem, Hodge
decomposition, the Hodge / de Rham rank computation, the Dolbeault
isomorphism, and Serre duality on a Riemann surface. Each one is a
\emph{major} classical-analysis result — none of them are a one-week
formalization, and Mathlib v4.28.0 contains none of them. This
chapter exists to make the full scope of those eight gaps visible: a
classical-mathematics walk-through, a Lean realisation plan, a
plain-English description, and an explicit list of sub-gaps that each
one in turn surfaces. The chapter is a roadmap, not a catalogue of
proofs. Every theorem in it is currently \emph{not ready}, and every
sub-leaf below the headline is a placeholder pinned to a Lean
declaration in \code{Jacobian/Analysis/} so that future workers
have a stable name to claim.
\end{layman}

The eight major analysis gaps in the Jacobian-challenge dep graph are
listed below.  Each gets its own subsection; each subsection contains
(i)~the headline classical statement, (ii)~a classical-proof outline,
(iii)~a Lean-formalization plan with the named sub-leaves, and
(iv)~a plain-English description for a non-specialist reader.  An
independent Lean build target lives at
\code{Jacobian/Analysis/<Node>/Overview.lean} for each node; sub-leaves
are pinned there as \code{theorem ... : True := trivial} placeholders
so the dependency graph has resolvable Lean targets even before any
substantive proof exists.

\begin{center}
\begin{tabular}{lll}
\textbf{Code} & \textbf{Headline} & \textbf{Build target} \\\hline
R1 & Radó's triangulation & \code{Jacobian.Analysis.Rado} \\
R2 & Tietze normal form & \code{Jacobian.Analysis.Tietze} \\
R3 & Polygonal-model theorem & \code{Jacobian.Analysis.PolygonalModel} \\
R4 & De Rham theorem & \code{Jacobian.Analysis.DeRham} \\
R5 & Hodge decomposition (Kähler) & \code{Jacobian.Analysis.HodgeDecomposition} \\
R6 & Hodge / de Rham rank & \code{Jacobian.Analysis.HodgeDeRham} \\
R7 & Dolbeault isomorphism & \code{Jacobian.Analysis.Dolbeault} \\
R8 & Serre duality on a Riemann surface & \code{Jacobian.Analysis.SerreDuality} \\
\end{tabular}
\end{center}

\bigskip

%======================================================================
\subsection{R1 — Radó's triangulation theorem}
\label{subsec:gap-R1-rado}

The headline (Theorem~\ref{input:rado-triangulation}) says: every
compact connected oriented topological 2-manifold admits a finite
triangulation.  This is the entry-point result of the entire
polygonal-model chain — without it, ``cut along a spanning tree to lay
the surface flat'' has no input.  Radó proved this in 1925; Doyle and
Moran (1968) gave the cleanest modern proof, simplified further by
Thomassen (1992).

\subsubsection*{Classical proof}

The proof has four phases:

\begin{enumerate}
\item \emph{Finite chart cover.}  By compactness, finitely many disk
  charts cover \(M\); each chart's image in \(\R^2\) contains an open
  Euclidean disk.
\item \emph{PL approximation.}  This is the dimension-2-specific
  step.  Any homeomorphism between open subsets of \(\R^2\) is
  uniformly approximable by a piecewise-linear (PL) homeomorphism
  (Doyle--Moran).  PL approximation can be carried out
  \emph{compatibly} across the finite atlas: at the cost of shrinking
  the atlas slightly, we get a cocycle of PL transition maps.  This
  is the step that fails in dimensions \(\geq 5\): in those
  dimensions, smoothable manifolds need not be triangulable
  (Kirby--Siebenmann).
\item \emph{Simplicial-complex assembly.}  Triangulate each PL chart
  by a fine simplicial subdivision of the disk image; take a common
  refinement across each overlap so adjacent triangulations agree
  edge-for-edge.  The result is a finite abstract simplicial
  complex.
\item \emph{Manifold link verification + homeomorphism.}  Verify the
  combinatorial 2-manifold conditions (every edge in exactly two
  triangles; every vertex link is a combinatorial circle).  The
  chart-by-chart maps glue into a continuous bijection from \(M\) to
  the geometric realisation; compact source + Hausdorff target +
  continuous bijection gives a homeomorphism.
\end{enumerate}

\subsubsection*{Lean proof plan}

Decomposed into 13 sub-leaves under
\code{Jacobian.Analysis.Rado.Overview}, organized in four phases
matching the classical proof.

\begin{itemize}
\item \emph{Phase 1 (already discharged).}
  \code{rado_finite_disk_atlas} (R1.1.1) is the only sub-leaf with a
  full Lean realisation: \code{compact_2manifold_has_finite_disk_atlas}
  in \code{StageA/RadoTheorem.lean}, sorry-free, ${\sim}40$ LOC using
  \code{IsCompact.elim_finite_subcover}.
\item \emph{Phase 2 (PL approximation).}  Sub-leaves R1.2.1--R1.2.4.
  Mathlib has no PL infrastructure; the entire phase needs to be
  built from scratch.  Estimated 600--900 LOC.
\item \emph{Phase 3 (simplicial-complex assembly).}  Sub-leaves
  R1.3.1--R1.3.4.  Depends on Phase 2.  Estimated 500--800 LOC plus
  the \code{IsCombinatorial2Manifold} predicate (R1-sub-C below).
\item \emph{Phase 4 (realisation homeomorphism).}  Sub-leaves
  R1.4.1--R1.4.3.  Phase 4's last step (R1.4.3) reduces to Mathlib's
  \code{Continuous.homeoOfEquivCompactToT2}; the rest is project-side
  glue.  Estimated 300--400 LOC.
\end{itemize}

\subsubsection*{Recursive sub-gaps surfaced}

\begin{itemize}
\item \textbf{R1-sub-A} (abstract simplicial complexes + geometric
  realisation): Mathlib has \code{SimplicialComplex} over a real
  vector space but \emph{not} the abstract combinatorial form with
  realisation as a topological quotient.  Sketch in
  \code{Jacobian/StageA/SimplicialComplex.lean}.  Estimated
  ~400 LOC.
\item \textbf{R1-sub-B} (Schoenflies-type theorem in dim 2 for PL
  approximation): a discrete topological theorem; classical and
  bounded, but Mathlib has nothing.  Estimated ~250 LOC.
\item \textbf{R1-sub-C} (\code{IsCombinatorial2Manifold} predicate):
  every edge in exactly two triangles, every vertex link a
  combinatorial circle.  Definition only ~30 LOC, but proving an
  example complex satisfies it requires assembled Phase~3.
\end{itemize}

\begin{theorem}[Gap R1: Radó's triangulation theorem (independent placeholder)]
\label{gap:R1-overview}
\uses{input:rado-triangulation}
The Lean target \code{Jacobian.Analysis.Rado.rado_overview} stubs
the headline; its sub-leaves R1.1.1--R1.4.3 are tracked in
\code{Jacobian.Analysis.Rado.Overview} and the bottom-up sketch
under \code{Jacobian/StageA/RadoTheorem.lean}.
\lean{JacobianChallenge.Analysis.Rado.rado_overview}
\notready
\end{theorem}

\begin{layman}
Radó's theorem says every compact two-dimensional surface — Riemann
surface or otherwise — can be cut into a finite collection of
triangles glued along their edges, in a way that respects the
topology exactly.  For a sphere, four triangles suffice; for a torus,
two; for higher genus, more, but always finitely many.  The classical
proof is a four-phase ballet: cover the surface with a finite atlas
of disks, smooth out the chart overlaps until they're piecewise
linear, fit the PL pieces together into a finite simplicial complex,
and verify the result really is a 2-manifold.  The \emph{piecewise
linear approximation} step is the secret ingredient, and it works
only in dimension \(\leq 3\).  In dimensions \(\geq 5\), there exist
manifolds with no triangulation at all (Kirby--Siebenmann's
counterexamples).  We're lucky that Riemann surfaces are
two-dimensional.

In Lean, the only piece of Radó that's currently sorry-free is the
finite chart cover (Phase~1), which is a one-line application of
compactness.  The other twelve sub-leaves are placeholders, and PL
approximation in particular is the largest item — Mathlib has no PL
infrastructure at all.  A full Lean proof of Radó alone is
estimated at 1500--2000 LOC, plus the abstract-simplicial-complex
package (R1-sub-A) of another ~400 LOC.
\end{layman}

%======================================================================
\subsection{R2 — Tietze normal form for orientable surface words}
\label{subsec:gap-R2-tietze}

The headline (Theorem~\ref{input:tietze-normal-form}) says: every edge
word arising from a polygonal presentation of a compact connected
orientable 2-manifold of genus \(g\) is Tietze-equivalent to the
standard relator
\[
  a_1\,b_1\,a_1^{-1}\,b_1^{-1}\,\cdots\,a_g\,b_g\,a_g^{-1}\,b_g^{-1}.
\]
This is the combinatorial heart of the topological-classification
theorem (Brahana 1921, Seifert--Threlfall 1934, often attributed in
modern textbooks to Massey).

\subsubsection*{Classical proof}

Reduction proceeds by induction on the length of the edge word:

\begin{enumerate}
\item \emph{Cyclic reduction.}  Use repeated inverse-pair cancellation
  to remove every \(\ell\,\ell^{-1}\) adjacency, including across the
  word's cyclic seam.  The result is a \emph{cyclically reduced}
  word.
\item \emph{Orientability constraint.}  In an orientable surface
  word, every generator letter appears exactly twice, once with each
  sign (so we never see \(a a\) or \(a^{-1} a^{-1}\) blocks).
\item \emph{Brahana handle creation.}  Find any pair
  \((\ell, \ell^{-1})\) and the pair \((m, m^{-1})\) interleaved
  with it.  By a sequence of cyclic rotations and
  letter-substitutions, ``slide'' the four letters together into a
  handle block \(\ell\,m\,\ell^{-1}\,m^{-1}\).
\item \emph{Strip and recurse.}  Once one handle block is at the
  start of the word, strip it off; the remainder is a shorter
  orientable cyclically reduced word, on which the induction
  hypothesis applies.  Concatenating the stripped handle gives the
  standard form.
\end{enumerate}

\subsubsection*{Lean proof plan}

Decomposed into 11 sub-leaves under
\code{Jacobian.Analysis.Tietze.Overview}, organized in four phases:

\begin{itemize}
\item \emph{Phase 1 (cyclic reduction).}  Sub-leaves R2.1.1--R2.1.2.
  Pure list-combinatorial; ~150 LOC of routine \code{List} work
  building on Mathlib's \code{List.IsRotated}.
\item \emph{Phase 2 (handle separation).}  Sub-leaves R2.2.1--R2.2.3.
  The Brahana handle-creation trick (R2.2.3) is the single deepest
  combinatorial step.  ~250 LOC.
\item \emph{Phase 3 (orientability flag preservation).}  Sub-leaves
  R2.3.1--R2.3.2.  Bookkeeping; ~100 LOC.
\item \emph{Phase 4 (induction on length).}  Sub-leaves
  R2.4.1--R2.4.3.  Standard structural induction; ~150 LOC.
\end{itemize}

\subsubsection*{Recursive sub-gaps surfaced}

\begin{itemize}
\item \textbf{R2-sub-A} (free-group surface presentation API on a
  finite alphabet).  Mathlib has \code{FreeGroup} but not the
  surface-presentation infrastructure; ~200 LOC.
\item \textbf{R2-sub-B} (cyclic-word equivalence on \code{List}).
  Mathlib has \code{List.IsRotated}; the cyclic-cancellation closure
  needs another ~80 LOC.
\item \textbf{R2-sub-C} (Brahana--Seifert--Threlfall handle-creation
  lemma, R2.2.3).  Pure combinatorics on \code{List Letter}; ~150
  LOC, but tricky.
\end{itemize}

\begin{theorem}[Gap R2: Tietze normal form (independent placeholder)]
\label{gap:R2-overview}
\uses{input:tietze-normal-form}
The Lean target \code{Jacobian.Analysis.Tietze.tietze_overview}
stubs the headline; its sub-leaves R2.1.1--R2.4.3 are tracked in
\code{Jacobian.Analysis.Tietze.Overview} alongside the bottom-up
sketch in \code{Jacobian/StageA/EdgeWordTietze.lean}.
\lean{JacobianChallenge.Analysis.Tietze.tietze_overview}
\notready
\end{theorem}

\begin{layman}
Pretend you've cut a doughnut surface open along loops and laid it
flat as a polygon.  Walking around the polygon's boundary, you see
each loop twice — once forward, once backward — so the boundary
traces out a sequence of letters like \(a b a^{-1} b^{-1}\).  Tietze
normal form says: regardless of which loops you happened to cut
along, the boundary word always simplifies, by a finite sequence of
``trivial'' moves, to the canonical pattern
\(a_1 b_1 a_1^{-1} b_1^{-1} \cdots a_g b_g a_g^{-1} b_g^{-1}\).  The
trivial moves are: cancel \(\ell \ell^{-1}\) (cutting along a loop
and immediately gluing it back is a no-op), and permute complete
handle blocks (relabeling which loop is first doesn't change the
surface).

The proof is an induction-on-length: cancel until cyclically reduced,
use orientability to ensure every letter appears with both signs,
slide one handle block to the front by the Brahana trick, strip it
off, recurse.  The Lean version mirrors this exactly, with the
deepest single move being the Brahana trick (R2.2.3).  Mathlib
has no group-presentation surface infrastructure, so all 11 sub-leaves
are currently placeholders; total estimated ~650 LOC for the
combinatorial side, plus ~280 LOC of supporting infrastructure
(R2-sub-A and R2-sub-B).
\end{layman}

%======================================================================
\subsection{R3 — The polygonal-model theorem}
\label{subsec:gap-R3-polygonal-model}

The headline (Theorem~\ref{thm:polygonal-model}) says: a compact
connected oriented Riemann surface of analytic genus \(g\) is
homeomorphic to the standard fundamental \(4g\)-gon.  This is the
\emph{assembly} gap: it consumes R1 (Radó) and R2 (Tietze) and outputs
the homeomorphism every period-lattice argument downstream relies on.
This is also the most-developed of the topology gaps; the project
already has substantial scaffolding under \code{Jacobian/Periods/}
and \code{Jacobian/StageA/}.

\subsubsection*{Classical proof}

\begin{enumerate}
\item \emph{Forget complex structure.}  A complex 1-manifold is
  automatically a smooth real 2-manifold, hence a topological
  2-manifold (already discharged sorry-free in Lean).
\item \emph{Apply Radó (R1).}  Get a finite triangulation of \(X\).
\item \emph{Cut along a spanning tree.}  Pick any spanning tree of
  the triangulation's 1-skeleton.  Cutting all non-tree edges
  produces a topological disk whose boundary is a closed walk of
  length \(2 \cdot (\text{number of non-tree edges})\) — and since
  the tree has \(V - 1\) edges out of \(V + F\) total, the
  Euler-formula constraint \(V - E + F = 2 - 2g\) makes the
  boundary walk's length exactly \(4g\).
\item \emph{Read off the edge word.}  Each non-tree edge produces
  two boundary appearances (one each side); orient by walking
  direction to get a four-letter alphabet \(a_i, b_i, a_i^{-1},
  b_i^{-1}\).  The boundary word is an arbitrary cyclically reduced
  word of length \(4g\) on this alphabet.
\item \emph{Apply Tietze (R2).}  Reduce the boundary word to the
  standard relator.
\item \emph{Lift Tietze moves to homeomorphisms.}  Each Tietze
  move (cancel an inverse pair; permute handle blocks) corresponds
  to a homeomorphism of polygons by cut-and-paste.  Compose all the
  moves to get a homeomorphism between the original cut-open polygon
  and the standard polygon.
\item \emph{Descend to the surface.}  Both polygons quotient to
  topologically the same surface (same gluing pattern post-Tietze);
  the homeomorphism descends to \(X \cong \mathrm{Polygon4g}(g)\).
\end{enumerate}

\subsubsection*{Lean proof plan}

Decomposed into 13 sub-leaves under
\code{Jacobian.Analysis.PolygonalModel.Overview}, plus 3 Stage-B
sub-leaves for the \(g_{\mathrm{an}} = g_{\mathrm{top}}\) bridge:

\begin{itemize}
\item \emph{Phase 1 (real-2-manifold structure).}  R3.1.1 is
  discharged sorry-free in
  \code{Jacobian/Periods/ComplexChartedSpaceReal2.lean}.  R3.1.2 is
  a pure invocation of R1.
\item \emph{Phase 2 (cut along spanning tree).}  R3.2.1 needs the
  spanning-tree existence theorem (Mathlib has
  \code{SimpleGraph.IsTree}; existence missing).  R3.2.2 and R3.2.3
  are routine combinatorics on the dual graph; project-side at
  \code{Jacobian/Periods/DualGraphCut.lean}.
\item \emph{Phase 3 (reduce the edge word).}  R3.3.1 invokes R2.
  R3.3.2 is the substantive geometric step: lifting Tietze moves to
  cut-and-paste homeomorphisms.  Estimated ~400 LOC.
\item \emph{Phase 4 (descend to surface).}  R3.4.1, R3.4.2: routine
  quotient-topology bookkeeping; ~150 LOC.
\item \emph{Stage-B bridge.}  R3.5.1 (smooth real 2-structure) is
  already done.  R3.5.2 (\(H_1\) of the polygon by cellular
  homology) and R3.5.3 (\(H_1\)-invariance under R3) are tracked at
  \code{Jacobian/Periods/H1EvenBasisViaSurfaceClassification.lean}
  and require \code{CellularHomologyRS} which is currently sorry.
\end{itemize}

\subsubsection*{Recursive sub-gaps surfaced}

\begin{itemize}
\item \textbf{R3-sub-A} (spanning-tree existence on a finite
  connected simple graph, R3.2.1).  ~80 LOC; sketched at
  \code{Jacobian/StageA/SpanningTree.lean}.
\item \textbf{R3-sub-B} (dual-graph + cut-along-tree combinatorics,
  R3.2.2 / R3.2.3).  ~250 LOC; tracked at
  \code{Jacobian/Periods/DualGraphCut.lean}.
\item \textbf{R3-sub-C} (Tietze-move-by-cut-and-paste correspondence,
  R3.3.2).  No Mathlib analogue; ~400 LOC of routine geometry but
  needs careful chart-management.
\end{itemize}

\begin{theorem}[Gap R3: Polygonal-model theorem (independent placeholder)]
\label{gap:R3-overview}
\uses{thm:polygonal-model,gap:R1-overview,gap:R2-overview}
The Lean target
\code{Jacobian.Analysis.PolygonalModel.polygonal_model_overview}
stubs the headline; sub-leaves R3.1.1--R3.5.3 are tracked in
\code{Jacobian.Analysis.PolygonalModel.Overview}.
\lean{JacobianChallenge.Analysis.PolygonalModel.polygonal_model_overview}
\notready
\end{theorem}

\begin{layman}
The polygonal model is the bridge between two pictures of a Riemann
surface.  Picture one: the curved surface in space, hard to
integrate over because it has holes around which Stokes' theorem
cannot apply.  Picture two: a flat polygon (a square for a torus, an
octagon for a genus-2 surface, a \(4g\)-gon for genus \(g\)), with
its boundary edges glued back together in the canonical pattern.
The polygon's interior is flat, simply connected, has no loops at
all — pedestrian for any classical-analysis argument.  All the
topological complexity has migrated to the boundary, where the
gluing pattern lives.

The proof is a geometric assembly: triangulate the surface (R1), cut
along a spanning tree to lay it flat as a polygon, read off the
boundary word, reduce that word to canonical form (R2).  Each step
in the reduction corresponds to a cut-and-paste homeomorphism of the
polygon, so the moves accumulate into a single homeomorphism between
the original cut-open polygon and the canonical \(4g\)-gon.  Descend
through the gluing quotient and you've identified the surface with
the canonical polygon.

In Lean, the polygon side of this story is already extensively
formalized — the standard \(4g\)-gon, edge words, cyclic reduction,
and the Tietze equivalence relation are all sorry-free.  What
remains sorry is the \emph{assembly}: getting the homeomorphism out
of R1 + R2 (especially R3.3.2, lifting Tietze moves to cut-and-paste
homeomorphisms).  Once R3 lands, the rest of the period-lattice
chain (Stokes on the polygon, Riemann bilinear, period-lattice
fullness) becomes accessible.
\end{layman}

%======================================================================
\subsection{R4 — The de Rham theorem}
\label{subsec:gap-R4-de-rham}

The headline (Theorem~\ref{input:de-rham-theorem}) says: for a smooth
manifold \(M\), integration of forms over chains gives an isomorphism
\(H^k_{\mathrm{dR}}(M, \C) \cong H^k_{\mathrm{sing}}(M, \C)\).  This is
de Rham's 1931 thesis result; the cleanest modern proofs go through
good covers (Bott--Tu) or sheaf cohomology.

\subsubsection*{Classical proof}

The Bott--Tu route runs as follows:

\begin{enumerate}
\item \emph{Differential-form package.}  Define \(\Omega^k(M)\) as
  smooth sections of \(\Lambda^k T^*M\), with exterior derivative
  \(d\), wedge \(\wedge\), and pullback \(f^*\).  Verify
  \(d^2 = 0\), Leibniz, and naturality.
\item \emph{Integration of forms.}  Define a smooth singular
  \(k\)-simplex \(\sigma : \Delta^k \to M\) (continuous, smooth in
  every chart).  Integration \(\int_\sigma \omega\) is well-defined
  on smooth simplices.  Stokes:
  \(\int_\sigma d\omega = \int_{\partial \sigma} \omega\).  The
  smooth-singular subcomplex of the singular chain complex is a
  quasi-isomorphism (smooth approximation of continuous chains).
\item \emph{Comparison map.}  Integration descends to
  \(H^k_{\mathrm{dR}} \to H^k_{\mathrm{sing}}\) by Stokes, and is
  natural and compatible with cup products.
\item \emph{Isomorphism via good covers.}  A \emph{good cover} is
  one where every finite intersection is contractible.  Both
  \(H^*_{\mathrm{dR}}\) and \(H^*_{\mathrm{sing}}\) restrict
  trivially to contractible opens, and both satisfy
  Mayer--Vietoris.  Induct on the size of a finite good cover
  using the five lemma; for general manifolds, take a colimit
  over all finite good subcovers.  Existence of good covers comes
  from a Riemannian metric and convex normal-coordinate balls.
\end{enumerate}

\subsubsection*{Lean proof plan}

Decomposed into 15 sub-leaves under
\code{Jacobian.Analysis.DeRham.Overview}, in four phases:

\begin{itemize}
\item \emph{Phase 1 (differential-form package).}  R4.1.1--R4.1.4.
  The largest single sub-gap; bundled \(\Omega^k\) needs to be built
  from Mathlib's \code{MFDeriv} primitives.  Estimated 800--1200
  LOC for a complete API.
\item \emph{Phase 2 (integration of forms).}  R4.2.1--R4.2.4.
  Smooth singular simplex (~50 LOC) + integration via measure on
  \(\Delta^k\) (~150 LOC) + Stokes for one simplex (~250 LOC) +
  smooth approximation (~200 LOC).
\item \emph{Phase 3 (cohomology comparison map).}  R4.3.1--R4.3.3.
  Routine assembly once Phases 1+2 are in hand; ~250 LOC.
\item \emph{Phase 4 (good-cover argument).}  R4.4.1--R4.4.4.
  Mayer--Vietoris is doable; the existence of good covers (R4.4.3)
  needs Riemannian-manifold infrastructure not in Mathlib.  ~600 LOC.
\end{itemize}

\subsubsection*{Recursive sub-gaps surfaced}

\begin{itemize}
\item \textbf{R4-sub-A} (bundled \(\Omega^k(M)\) as a
  \(C^\infty(M)\)-module, R4.1).  ~400 LOC; foundational for
  R5/R6/R7 also.
\item \textbf{R4-sub-B} (smooth singular chains and Stokes on a
  smooth \(k\)-simplex, R4.2.3).  ~300 LOC building on Mathlib's
  measure-theoretic Stokes for boxes.
\item \textbf{R4-sub-C} (good-cover existence on a smooth manifold,
  R4.4.3).  ~250 LOC plus prerequisite Riemannian-metric package
  (~500 LOC of foundational work — see R5).
\end{itemize}

\begin{theorem}[Gap R4: De Rham theorem (independent placeholder)]
\label{gap:R4-overview}
\uses{input:de-rham-theorem}
The Lean target \code{Jacobian.Analysis.DeRham.deRham_overview}
stubs the headline; sub-leaves R4.1.1--R4.4.4 are tracked in
\code{Jacobian.Analysis.DeRham.Overview} and the bottom-up sketch
in \code{Jacobian/StageB/DeRhamComparison.lean}.
\lean{JacobianChallenge.Analysis.DeRham.deRham_overview}
\notready
\end{theorem}

\begin{layman}
There are two ways to count holes in a smooth manifold.  The
\emph{differential-forms} way: a closed form is one with
\(d\omega = 0\); an exact form is one of the shape \(\omega = d\alpha\);
the quotient (closed mod exact) is the de Rham cohomology.  The
\emph{singular-chains} way: a cycle is a formal combination of
simplices with \(\partial = 0\); a boundary is one of the shape
\(\partial \beta\); the dual quotient is the singular cohomology.
De Rham's theorem says these two pictures agree, and the
isomorphism is integration: pair a closed form against a cycle and
read off a number.

Why bother?  The de Rham side is the natural language of analysis on
manifolds (everything is a form, integration is automatic).  The
singular side is the natural language of topology (cycles are loops,
boundaries are filled-in regions).  De Rham unites them.  Without
de Rham, there's no way to compare the analytic genus
(``how many holomorphic 1-forms?'') with the topological genus
(``how many independent loops?'').

The Lean version of this proof is large, partly because Mathlib has
no bundled differential-form package at all — the
\(\Omega^k\)-module of smooth \(k\)-forms with \(d\), \(\wedge\),
pullback all simultaneously linear and natural, has to be built from
scratch (R4-sub-A).  The good-cover argument also pulls in
Riemannian-metric infrastructure (R4-sub-C, ~500 LOC of
prerequisite).  Total estimated for R4 + supporting sub-gaps:
2500--3500 LOC.
\end{layman}

%======================================================================
\subsection{R5 — Hodge decomposition on a compact Kähler manifold}
\label{subsec:gap-R5-hodge}

The headline (Theorem~\ref{input:hodge-decomposition}) says: on a
compact Kähler manifold (in particular every compact Riemann surface
with any Hermitian metric), the Laplace--Beltrami operator
\(\Delta = d\delta + \delta d\) has finite-dimensional kernel, every
de Rham cohomology class has a unique harmonic representative, and on
a Kähler manifold harmonic forms split by \((p,q)\)-bidegree.

\subsubsection*{Classical proof}

This is the deepest of the eight gaps and the proof has the most
moving pieces:

\begin{enumerate}
\item \emph{Riemannian metric and Hodge \(*\).}  Pick a
  Riemannian metric \(g\); it induces a metric on cotangent space and
  hence on \(\Lambda^k T^* M\), plus a volume form.  The Hodge star
  \(* : \Omega^k \to \Omega^{n-k}\) is the fibrewise duality.
\item \emph{Codifferential and Laplacian.}  Define \(\delta = (-1)^* * d *\)
  (formal \(L^2\)-adjoint of \(d\)) and \(\Delta = d\delta + \delta d\).
  \(\Delta\) is symmetric, non-negative, commutes with \(d\),
  \(\delta\), \(*\).
\item \emph{Ellipticity + elliptic regularity.}  \(\Delta\) is an
  elliptic second-order differential operator.  Its principal
  symbol is \(|\xi|^2 \cdot \mathrm{id}\), invertible for
  \(\xi \neq 0\).  Elliptic regularity: every distributional
  solution of \(\Delta \omega = f\) with smooth \(f\) is itself
  smooth.  Combined with Rellich--Kondrachov compactness on a
  compact manifold, this makes \(\Delta : H^s_{k+2} \to H^s_k\) (in
  Sobolev norms on forms) Fredholm.
\item \emph{Harmonic decomposition.}  The Fredholm property gives
  \(\Omega^k(M) = \mathcal{H}^k(M) \oplus d(\Omega^{k-1}) \oplus \delta(\Omega^{k+1})\),
  an \(L^2\)-orthogonal decomposition, with
  \(\mathcal{H}^k = \ker \Delta\) the (finite-dimensional) harmonic
  forms.  Every de Rham class has a unique harmonic representative.
\item \emph{Kähler bigrading.}  On a Kähler manifold, the metric,
  complex structure, and Kähler form are mutually compatible.  The
  Kähler identities \([\Lambda, \partial] = i\bar\partial^*\) (and
  cousins) imply \(\Delta_d = 2\Delta_{\bar\partial} = 2\Delta_\partial\).
  Hence \(\Delta\) preserves \((p,q)\)-bidegree, the harmonic
  decomposition splits as
  \(\mathcal{H}^k = \bigoplus_{p+q=k} \mathcal{H}^{p,q}\), and Hodge
  symmetry \(\mathcal{H}^{p,q} = \overline{\mathcal{H}^{q,p}}\)
  follows from complex conjugation.
\end{enumerate}

\subsubsection*{Lean proof plan}

Decomposed into 19 sub-leaves under
\code{Jacobian.Analysis.HodgeDecomposition.Overview}, in six phases.
The largest single piece is the Sobolev-elliptic-regularity package
(R5-sub-B); without that, the entire decomposition argument has no
foundation.

\begin{itemize}
\item \emph{Phase 1 (Riemannian + Hodge \(*\)).}  R5.1.1--R5.1.4.
  Needs Riemannian-metric infrastructure (R5-sub-A).  ~400 LOC.
\item \emph{Phase 2 (codifferential + Laplacian).}  R5.2.1--R5.2.3.
  ~200 LOC once Phase 1 lands.
\item \emph{Phase 3 (ellipticity + regularity).}  R5.3.1--R5.3.4.
  Mathlib has neither distributional Sobolev spaces on a manifold,
  nor elliptic regularity, nor Rellich--Kondrachov for forms.
  Estimated 1500--2500 LOC for the package; the largest single
  obstacle to formalizing R5.
\item \emph{Phase 4 (harmonic decomposition).}  R5.4.1--R5.4.4.
  Routine functional analysis once Phase 3 is in hand; ~300 LOC.
\item \emph{Phase 5 (Kähler bigrading).}  R5.5.1--R5.5.5.  Mostly
  algebraic; ~400 LOC.
\item \emph{Phase 6 (Riemann surface specialisation).}  R5.6.1,
  R5.6.2.  ~150 LOC.
\end{itemize}

\subsubsection*{Recursive sub-gaps surfaced}

\begin{itemize}
\item \textbf{R5-sub-A} (Riemannian metric on a manifold).
  Mathlib has none; needs bundle structure on \(T^*M\).  ~500 LOC,
  shared with R4-sub-C.
\item \textbf{R5-sub-B} (Sobolev spaces of forms + elliptic regularity
  for second-order operators).  The single largest sub-gap in the
  entire eight-node tree; estimated 2500--4000 LOC.  Mathlib has
  Sobolev spaces on \(\R^n\) (\code{Mathlib.Analysis.Distribution})
  but not on manifolds, and not for vector-bundle-valued forms.
\item \textbf{R5-sub-C} (Rellich--Kondrachov compact embedding for
  forms).  ~400 LOC after R5-sub-B.
\item \textbf{R5-sub-D} (Fredholm-alternative on a compact manifold).
  ~300 LOC after R5-sub-B + R5-sub-C.
\item \textbf{R5-sub-E} (Kähler-identity package).  ~400 LOC of
  algebraic-geometry-flavored pencils; needs Phase 5 sub-leaves
  R5.5.1--R5.5.5.
\end{itemize}

\begin{theorem}[Gap R5: Hodge decomposition (independent placeholder)]
\label{gap:R5-overview}
\uses{input:hodge-decomposition}
The Lean target
\code{Jacobian.Analysis.HodgeDecomposition.hodge_decomposition_overview}
stubs the headline; sub-leaves R5.1.1--R5.6.2 are tracked in
\code{Jacobian.Analysis.HodgeDecomposition.Overview}.
\lean{JacobianChallenge.Analysis.HodgeDecomposition.hodge_decomposition_overview}
\notready
\end{theorem}

\begin{layman}
The Hodge decomposition theorem is the most analytic of the eight
gaps and the one that pulls in the largest piece of missing
foundational work in Mathlib.  The headline statement is:
on a compact Kähler manifold (a Riemann surface with any choice of
Hermitian metric counts) every cohomology class has a \emph{unique}
harmonic representative — a form lying in the kernel of the
Laplace--Beltrami operator \(\Delta = d\delta + \delta d\).  On a
compact Riemann surface, harmonic 1-forms further split by type:
holomorphic 1-forms (\((1,0)\)-type) and antiholomorphic 1-forms
(\((0,1)\)-type), each \(g\)-dimensional, summing to \(2g\).

The proof's deepest layer is elliptic-PDE machinery: ellipticity of
\(\Delta\) makes it (modulo a finite-dimensional null space) an
isomorphism on appropriately defined Sobolev spaces of forms.
That's what gives unique harmonic representatives.  The Kähler
identities then sharpen the picture by revealing that the Laplacian
respects the \((p,q)\)-bidegree decomposition, so harmonic forms
themselves split by type.

For Lean, the sticking point is that Mathlib v4.28.0 has no Sobolev
spaces of forms on a manifold and no elliptic regularity for
manifold-level operators.  Building this from scratch is the
single largest task in the eight-gap tree (R5-sub-B alone:
2500--4000 LOC).  Once that's in hand, the rest of the decomposition
(harmonic representatives, Kähler bigrading, Hodge symmetry) is
routine modern functional analysis.  Total estimated for R5 with all
sub-gaps: 5500--8500 LOC.
\end{layman}

%======================================================================
\subsection{R6 — Hodge / de Rham rank for a Riemann surface}
\label{subsec:gap-R6-hodge-deRham}

The headline (Theorem~\ref{input:hodge-deRham}) says: for a compact
connected oriented Riemann surface \(X\) of analytic genus \(g\),
\(\mathrm{rank}_{\Z}\,H_1(X, \Z) = 2g\).  This is the assembly
gap that consumes R4 (de Rham), R5 (Hodge decomposition), and R7
(Dolbeault) and outputs the integer rank used by the
period-lattice fullness argument.

\subsubsection*{Classical proof}

A three-step assembly:

\begin{enumerate}
\item \emph{Hodge \(\Rightarrow\)
  \(\dim_{\C} H^1_{\mathrm{dR}}(X, \C) = 2g\).}  Apply R5 to split
  \(H^1_{\mathrm{dR}}(X, \C) = H^{1,0}(X) \oplus H^{0,1}(X)\).
  Apply R7 to identify \(H^{1,0}(X) \cong H^0(X, \Omega^1_X)\), which
  by definition has dimension \(g\).  Apply Hodge symmetry
  (R5.5.5) to get \(\dim H^{0,1} = g\).
\item \emph{De Rham \(\Rightarrow\)
  \(H^1_{\mathrm{sing}}(X, \R) \cong H^1_{\mathrm{dR}}(X, \R)\).}
  Apply R4.  Real dimension is \(2g\).
\item \emph{Universal coefficients \(\Rightarrow\)
  \(\mathrm{rank}_{\Z} H_1(X, \Z) = \dim_\R H^1_{\mathrm{sing}}(X, \R)\).}
  In degree 1 on a compact CW manifold, the universal-coefficient
  theorem gives \(H^1(X, \R) \cong \mathrm{Hom}(H_1(X, \Z), \R)\),
  whose real dimension equals the integer rank of \(H_1(X, \Z)\).
\end{enumerate}

\subsubsection*{Lean proof plan}

Decomposed into 11 sub-leaves under
\code{Jacobian.Analysis.HodgeDeRham.Overview}, in four phases:

\begin{itemize}
\item \emph{Phase 1 (real dim of \(H^1_{\mathrm{dR}}\)).}
  R6.1.1--R6.1.5.  All routine assembly above R5/R7; ~200 LOC.
\item \emph{Phase 2 (de Rham comparison).}  R6.2.1, R6.2.2.  Pure
  invocation of R4; ~50 LOC.
\item \emph{Phase 3 (singular-to-integer-homology).}  R6.3.1--R6.3.4.
  Universal coefficients and rank-Hom lemmas; Mathlib has the algebra
  but not the singular variant for manifolds.  ~400 LOC.
\item \emph{Phase 4 (final).}  R6.4.1.  ~30 LOC.
\end{itemize}

R6 is small in itself (estimated ~700 LOC) but it depends on three of
the eight major gaps simultaneously; it cannot be discharged in
isolation.

\subsubsection*{Recursive sub-gaps surfaced}

\begin{itemize}
\item \textbf{R6-sub-A} (universal-coefficient theorem in degree 1
  for a CW manifold, R6.3.1).  Algebraic UCT exists in Mathlib
  (\code{Mathlib.Algebra.Homology.UniversalCoefficient});
  geometric variant for \(H^*_{\mathrm{sing}}(X, \Z)\) does not.
  ~250 LOC.
\item \textbf{R6-sub-B} (\(H_1(X, \Z)\) finitely generated for a
  compact CW manifold, R6.3.2).  Mathlib has \(H_*\) for chain
  complexes but not the manifold-level finite-CW theorem.
  ~150 LOC, after a triangulation is in hand (so this depends on R1).
\item \textbf{R6-sub-C} (\(\dim_\R \mathrm{Hom}(A, \R) = \mathrm{rank}_\Z A\)
  for a fg abelian group, R6.3.3).  Routine; ~80 LOC; Mathlib has
  the building blocks.
\end{itemize}

\begin{theorem}[Gap R6: Hodge / de Rham rank (independent placeholder)]
\label{gap:R6-overview}
\uses{input:hodge-deRham,gap:R4-overview,gap:R5-overview,gap:R7-overview}
The Lean target \code{Jacobian.Analysis.HodgeDeRham.hodge_deRham_overview}
stubs the headline; sub-leaves R6.1.1--R6.4.1 are tracked in
\code{Jacobian.Analysis.HodgeDeRham.Overview}.
\lean{JacobianChallenge.Analysis.HodgeDeRham.hodge_deRham_overview}
\notready
\end{theorem}

\begin{layman}
R6 is the simplest of the eight in terms of its own work, but it
sits at the confluence of three other gaps — it can't be done
without R4, R5, and R7 all in place.  The headline says: on a genus-\(g\)
Riemann surface, the integer first homology group has rank exactly
\(2g\).  In words: there are \(2g\) topologically independent loops
on \(X\), no more, no less.

The proof is a chain of three identifications.  First, Hodge (R5) +
Dolbeault (R7) say that the complex de Rham cohomology has dimension
\(2g\) — \(g\) holomorphic 1-forms plus \(g\) antiholomorphic ones.
Second, de Rham (R4) says singular cohomology over \(\R\) matches
de Rham cohomology, so the real singular cohomology also has
dimension \(2g\).  Third, the universal-coefficient theorem
identifies the real cohomology dimension with the integer homology
rank, giving \(\mathrm{rank}_\Z H_1(X, \Z) = 2g\).

Why bother?  This is the rank fact every period-lattice argument in
the project relies on: the symplectic basis \(a_1, b_1, \ldots, a_g,
b_g\) generates \(H_1(X, \Z)\), so the period vectors compute a full
rank-\(2g\) lattice in \((H^0(X, \Omega^1))^\vee \cong \C^g\).  No
fullness, no Jacobian.  R6 is short in itself (~700 LOC) but
unblockable until R4 + R5 + R7 are in hand, which is the dependency
shape that makes the whole tree expensive.
\end{layman}

%======================================================================
\subsection{R7 — The Dolbeault isomorphism}
\label{subsec:gap-R7-dolbeault}

The headline (Theorem~\ref{input:dolbeault}) says: for a complex
manifold \(X\) and a holomorphic vector bundle \(E\),
\(H^q_{\bar\partial}(X, E) \cong H^q(X, \mathcal{O}(E))\) — Dolbeault
cohomology agrees with sheaf cohomology.  On a Riemann surface with
\(E = \mathcal{O}_X\): \(H^{0,1}_{\bar\partial}(X) \cong H^1(X, \mathcal{O}_X)\).
This is the analytic--algebraic bridge that lets Hodge-theoretic
content (\(L^2\) inner products, harmonic representatives) be moved
to sheaf-cohomological content (the natural language of Riemann--Roch
and Serre duality).

\subsubsection*{Classical proof}

\begin{enumerate}
\item \emph{Dolbeault complex.}  On a complex manifold, forms split
  by bidegree, \(\Omega^k(X) = \bigoplus_{p+q=k} \Omega^{p,q}(X)\),
  and \(d\) splits as \(\partial + \bar\partial\).  Define
  Dolbeault cohomology \(H^{p,q}_{\bar\partial}(X) :=
  \ker(\bar\partial : \Omega^{p,q} \to \Omega^{p,q+1}) /
  \mathrm{im}(\bar\partial : \Omega^{p,q-1} \to \Omega^{p,q})\).
\item \emph{\(\bar\partial\)-Poincaré on a polydisk.}  On a polydisk
  \(U \subseteq \C^n\), every \(\bar\partial\)-closed
  \((p,q)\)-form with \(q \geq 1\) is locally
  \(\bar\partial\)-exact.  In one variable: solve
  \(\bar\partial f = g\) by the Cauchy--Pompeiu integral
  \(f(z) = \frac{1}{2\pi i} \int_U g(w)/(w - z) \, dw \wedge d\bar w\).
\item \emph{Fine resolution of \(\Omega^p_X\).}  The sheafified
  Dolbeault complex
  \(0 \to \Omega^p_X \to \Omega^{p,0} \to \Omega^{p,1} \to \cdots\)
  is exact (by the Poincaré lemma) and each \(\Omega^{p,q}\) is a
  fine sheaf (it has smooth partitions of unity).
\item \emph{Sheaf cohomology via fine resolution.}  Sheaf cohomology
  can be computed from any fine resolution; therefore
  \(H^q(X, \Omega^p_X) \cong H^q_{\bar\partial}(X)\).  Same proof
  works for \(E\)-valued forms with \(E\) holomorphic.
\end{enumerate}

\subsubsection*{Lean proof plan}

Decomposed into 12 sub-leaves under
\code{Jacobian.Analysis.Dolbeault.Overview}, in four phases:

\begin{itemize}
\item \emph{Phase 1 (Dolbeault complex).}  R7.1.1--R7.1.4.  Bigraded
  forms (R7-sub-A) — Mathlib has nothing.  ~600 LOC.
\item \emph{Phase 2 (\(\bar\partial\)-Poincaré).}  R7.2.1, R7.2.2.
  The one-dimensional case reduces to the Cauchy--Pompeiu integral,
  which Mathlib has via \code{Complex.integral_cauchy_pompeiu}.
  ~250 LOC.
\item \emph{Phase 3 (sheaf comparison).}  R7.3.1--R7.3.3.  Needs
  derived-functor / acyclic-resolution machinery on the abelian
  category of sheaves.  Mathlib has Čech cohomology
  (\code{Mathlib.Topology.Sheaves.Cech}); the comparison-via-fine-resolution
  needs explicit acyclicity.  ~400 LOC.
\item \emph{Phase 4 (Riemann surface specialisation).}
  R7.4.1--R7.4.3.  ~150 LOC.
\end{itemize}

\subsubsection*{Recursive sub-gaps surfaced}

\begin{itemize}
\item \textbf{R7-sub-A} (bigraded forms \(\Omega^{p,q}\) on a complex
  manifold).  Mathlib has neither bigraded forms nor \(\partial\) /
  \(\bar\partial\).  ~600 LOC, partially shared with R5.
\item \textbf{R7-sub-B} (\(\bar\partial\)-Poincaré on a polydisk).
  The one-variable case is ~80 LOC building on Cauchy--Pompeiu.
  Multivariable case is induction; ~250 LOC.
\item \textbf{R7-sub-C} (fine sheaves on a smooth manifold,
  R7.2.2).  Needs the smooth-partition-of-unity API for sheaves.
  ~250 LOC.
\item \textbf{R7-sub-D} (sheaf cohomology agrees with cohomology of a
  fine resolution, R7.3.1).  Standard derived-functor argument; ~300
  LOC.
\end{itemize}

\begin{theorem}[Gap R7: Dolbeault isomorphism (independent placeholder)]
\label{gap:R7-overview}
\uses{input:dolbeault}
The Lean target \code{Jacobian.Analysis.Dolbeault.dolbeault_overview}
stubs the headline; sub-leaves R7.1.1--R7.4.3 are tracked in
\code{Jacobian.Analysis.Dolbeault.Overview}.
\lean{JacobianChallenge.Analysis.Dolbeault.dolbeault_overview}
\notready
\end{theorem}

\begin{layman}
On a complex manifold, differential forms split by bidegree:
a 1-form is either of type \((1,0)\) — looking like \(f \, dz\) — or
of type \((0,1)\) — looking like \(g \, d\bar z\) — or a sum of both.
The \(\bar\partial\) operator sends \((p,q)\)-forms to
\((p,q+1)\)-forms; its cohomology \(H^{p,q}_{\bar\partial}\) is the
analytic side of Dolbeault.

The algebraic side, sheaf cohomology \(H^q(X, \Omega^p_X)\), is what
Riemann--Roch and Serre duality naturally speak.  Dolbeault's theorem
says these two pictures agree.  On a Riemann surface and \(p=q=0\):
\(H^0\) is global functions, in both pictures the same.  On a Riemann
surface and \((p,q) = (0,1)\): \(H^{0,1}_{\bar\partial}\) is
``\(\bar\partial\)-cohomology of \(\bar z\)-forms,'' and Dolbeault
identifies it with the structure-sheaf cohomology
\(H^1(X, \mathcal{O}_X)\) — exactly the sheaf-side input to Serre
duality.

The proof proceeds by showing the Dolbeault complex is a fine
resolution of the holomorphic-form sheaf, hence computes its
cohomology.  The fine-resolution side is locally trivial — the key
ingredient is the \(\bar\partial\)-Poincaré lemma, which on a
polydisk reduces to the Cauchy--Pompeiu integral formula.  Mathlib
has the integral formula, but no bigraded forms (R7-sub-A is shared
with R5) and no fine-sheaf API on manifolds.  Total estimated for R7:
~1300 LOC plus shared infrastructure.
\end{layman}

%======================================================================
\subsection{R8 — Serre duality on a compact Riemann surface}
\label{subsec:gap-R8-serre}

The headline (Theorem~\ref{thm:serre-duality-rs}) says: for a compact
Riemann surface \(X\) and a coherent sheaf \(F\), the cup-product /
trace pairing \(H^q(X, F) \times H^{1-q}(X, \mathcal{H}om(F, K_X))
\to H^1(X, K_X) \to \C\) is non-degenerate.  In the case \(F =
\mathcal{O}_X\), \(q = 1\): \(H^1(X, \mathcal{O}_X) \cong
H^0(X, K_X)^\vee\), which forces \(\dim H^1(X, \mathcal{O}_X) = g
= \dim H^0(X, \Omega^1)\).

\subsubsection*{Classical proof}

\begin{enumerate}
\item \emph{Dualizing sheaf.}  Identify \(K_X = \Omega^1_X\) as the
  dualizing sheaf — a line bundle.  Define \(\mathcal{H}om(F, K_X)\),
  the canonical Serre dual of any coherent \(F\).
\item \emph{Trace map.}  Define a residue map at each point of \(X\)
  on meromorphic 1-forms; aggregate via Čech representatives to a
  trace \(\mathrm{tr} : H^1(X, K_X) \to \C\).  The classical
  \emph{residue theorem} (sum of residues of a global meromorphic
  1-form is zero) shows \(\mathrm{tr}\) is well-defined and is an
  isomorphism.
\item \emph{Cup-product pairing.}  Cup product
  \(\smile : H^q(X, F) \times H^{1-q}(X, \mathcal{H}om(F, K_X)) \to
  H^1(X, F \otimes \mathcal{H}om(F, K_X))\), composed with the
  evaluation \(F \otimes \mathcal{H}om(F, K_X) \to K_X\), gives a
  pairing into \(H^1(X, K_X)\); composing with \(\mathrm{tr}\) gives
  the Serre pairing.
\item \emph{Non-degeneracy via harmonic forms.}  Identify both
  sides of the pairing with harmonic-form spaces (using R5 + R7).
  The pairing realises the \(L^2\) inner product on harmonic forms,
  which is positive-definite, hence non-degenerate.
\end{enumerate}

\subsubsection*{Lean proof plan}

R8 is the most-developed of the eight gaps: the project already has
~32 files under \code{Jacobian/HolomorphicForms/Serre/} that
decompose \code{serre_duality_rs} into ~80 named sub-leaves.  This
section's umbrella decomposes those further into 13 phase-named
sub-leaves under \code{Jacobian.Analysis.SerreDuality.Overview}, in
five phases:

\begin{itemize}
\item \emph{Phase 1 (dualizing sheaf).}  R8.1.1--R8.1.3.  Mostly
  done at the Lean stub level: \code{RSDualizingSheaf} exists,
  \code{Ω¹_X} as a line bundle is sketched.  ~250 LOC remaining.
\item \emph{Phase 2 (trace map).}  R8.2.1--R8.2.3.  Needs the
  residue theorem (R8-sub-C), which is classical complex analysis on
  a Riemann surface.  ~400 LOC.
\item \emph{Phase 3 (cup-product pairing).}  R8.3.1--R8.3.3.  Mathlib
  has the cup product on simplicial cohomology; Čech variant is
  R8-sub-B.  ~350 LOC.
\item \emph{Phase 4 (non-degeneracy).}  R8.4.1, R8.4.2.  This is the
  R5/R7-dependent step.  ~300 LOC modulo R5 + R7.
\item \emph{Phase 5 (Riemann surface specialisation).}
  R8.5.1--R8.5.3.  ~150 LOC.
\end{itemize}

The existing 32-file Serre tree exposes much finer-grained sub-leaves
than the 13 listed here; the numbering above is the umbrella view.

\subsubsection*{Recursive sub-gaps surfaced}

\begin{itemize}
\item \textbf{R8-sub-A} (sheaf cohomology of an analytic line bundle
  on a Riemann surface).  Sketched at
  \code{Jacobian/HolomorphicForms/SheafCohomologyRS.lean}.
  ~600 LOC.
\item \textbf{R8-sub-B} (cup product on Čech cohomology).  Mathlib
  has the cup product on simplicial cohomology and partial Čech
  apparatus; the cup product on Čech needs porting.  ~300 LOC.
\item \textbf{R8-sub-C} (residue theorem on a compact Riemann
  surface, R8.2.3).  Classical: every global meromorphic 1-form has
  residues summing to zero.  ~250 LOC; assumes Stokes on a Riemann
  surface with boundary (a sibling project gap).
\item \textbf{R8-sub-D} (\(L^2\)-realisation of the Serre pairing
  via harmonic forms, R8.4.2).  Bridges R5 + R7 + R8.  ~200 LOC.
\end{itemize}

\begin{theorem}[Gap R8: Serre duality (independent placeholder)]
\label{gap:R8-overview}
\uses{thm:serre-duality-rs,gap:R5-overview,gap:R7-overview}
The Lean target
\code{Jacobian.Analysis.SerreDuality.serre_duality_overview} stubs
the headline; sub-leaves R8.1.1--R8.5.3 are tracked in
\code{Jacobian.Analysis.SerreDuality.Overview} and the 32-file
Serre tree at \code{Jacobian/HolomorphicForms/Serre/}.
\lean{JacobianChallenge.Analysis.SerreDuality.serre_duality_overview}
\notready
\end{theorem}

\begin{layman}
Serre duality is a magical identity in complex algebraic geometry.
On a compact Riemann surface, the first sheaf cohomology of the
structure sheaf, \(H^1(X, \mathcal{O}_X)\), is the dual vector
space to the global sections of the canonical bundle,
\(H^0(X, K_X) = \) holomorphic 1-forms.  Dual: not just isomorphic
in dimension, but explicitly paired by an integration identity on
\(X\).

In particular: \(\dim H^1(X, \mathcal{O}_X) = \dim H^0(X, K_X) = g\)
— the genus.  This is the key ingredient in Riemann--Roch (which
otherwise has only the inequality \(h^0 - h^1 \geq d - g + 1\); Serre
upgrades it to an equality).  It's also what lets us identify the
analytic Jacobian (\(H^0(X, \Omega^1)^\vee / H_1(X, \Z)\)) with the
algebraic Jacobian (\(\mathrm{Pic}^0(X)\)) — the deepest cross-bridge
in the project.

The proof has four moving parts.  First, define a trace map from the
top sheaf cohomology to \(\C\), via residues on a Čech cover.
Second, define a cup-product pairing using the cup product on Čech
cohomology and the natural evaluation map on sheaves.  Third,
identify both sides of the pairing with harmonic-form spaces using
R5 + R7.  Fourth, observe that the pairing realises the
\(L^2\) inner product on harmonic forms, which is positive-definite
hence non-degenerate.

R8 is the most-developed Lean gap of the eight: ~32 files already
sketch the proof tree.  But many leaves remain sorry, and several
(the \(L^2\)-bridge to R5 + R7, the residue theorem) are themselves
non-trivial sub-gaps.  Total estimated for R8 alone: ~3000 LOC
modulo R5 + R7.  The dependency on R5 is what makes R8 expensive —
R8 doesn't need full Hodge decomposition, but it does need the
\(L^2\)-pairing-positive-definiteness ingredient that comes from the
Hodge package.
\end{layman}

%======================================================================
\subsection{Summary of recursive sub-gaps}
\label{subsec:gap-summary}

The eight-node tree, when expanded one level into recursive sub-gaps,
surfaces a roughly 25-node graph.  Several sub-gaps are shared
across multiple parent nodes (the bundled \(\Omega^k\) package
appears under R4, R5, and R7; the Riemannian metric appears under R4
and R5; the abstract simplicial complex appears under R1 and R3).
A consolidated estimate of the total Lean-formalisation cost,
including all shared sub-gaps counted once, is:

\begin{center}
\begin{tabular}{lr}
\textbf{Component} & \textbf{Estimated LOC} \\\hline
R1 + R1-sub-\{A,B,C\} & 2000--2700 \\
R2 + R2-sub-\{A,B,C\} & 900--1100 \\
R3 + R3-sub-\{A,B,C\} & 1500--2000 \\
R4 + R4-sub-\{A,B,C\} & 2500--3500 \\
R5 + R5-sub-\{A,B,C,D,E\} & 5500--8500 \\
R6 + R6-sub-\{A,B,C\} & 1100--1400 \\
R7 + R7-sub-\{A,B,C,D\} & 2200--2900 \\
R8 + R8-sub-\{A,B,C,D\} & 3300--4300 \\\hline
\textbf{Total (with shared deduplication)} & \textbf{15000--22000 LOC} \\
\end{tabular}
\end{center}

The dominant single sub-gap is R5-sub-B (Sobolev / elliptic
regularity for forms on a manifold), which alone is estimated at
2500--4000 LOC and is on the critical path for both R5 and R7.

\bigskip

\begin{layman}
Adding it all up: the eight major analysis gaps, expanded one level
into their internal sub-gaps, total around fifteen to twenty-two
thousand lines of Lean.  That's a substantial fraction of the entire
project but it's a knowable one — every node is named, every
dependency is identified, every sub-gap has at least a placeholder
Lean target.  The single biggest cost driver is \emph{Sobolev spaces
of forms on a manifold plus elliptic regularity} (R5-sub-B), which
underpins the Hodge decomposition (R5) and indirectly the Dolbeault
proof (R7).  Once that's in hand, much of the rest cascades.

The other large cost is \emph{bundled differential forms}
(R4-sub-A): a \(C^\infty(M)\)-module \(\Omega^k(M)\) with all the
expected operators (\(d\), \(\wedge\), pullback) simultaneously
linear, natural, and compatible.  Mathlib has the pieces
(\code{MFDeriv}, exterior algebra in finite dimensions, manifold
charts) but no integrated bundled form package; this is shared
infrastructure across R4, R5, and R7.

Among the eight, R6 is the smallest in itself but the most
\emph{blocked}: it depends on R4, R5, and R7 simultaneously, so it
cannot be touched until all three are deep enough to expose their
output statements.  R3, by contrast, is the most \emph{ready}: its
Lean infrastructure is already largely in place, blocked mainly on
R1 (Radó) and R2 (Tietze).  A pragmatic first move would be to
formalise R1 + R2 + R3 as a topology-only sub-project, in parallel
with R4-sub-A (bundled \(\Omega^k\)) and R5-sub-A (Riemannian
metric) as foundational analytic infrastructure.  Both of those
unblock the rest of the analytic chain.
\end{layman}
